\documentclass[11pt, a4paper]{article}

% ── Packages ──────────────────────────────────────────────────────────
\usepackage[utf8]{inputenc}
\usepackage[T1]{fontenc}
\usepackage{lmodern}
\usepackage[margin=1in]{geometry}
\usepackage{amsmath, amssymb, amsthm}
\usepackage{booktabs}
\usepackage{longtable}
\usepackage{graphicx}
\usepackage{hyperref}
\usepackage{xcolor}
\usepackage{listings}
\usepackage{tcolorbox}
\usepackage{polya}

\hypersetup{
  colorlinks=true,
  linkcolor=polyablue,
  urlcolor=polyablue,
  citecolor=polyablue,
}

% ── Code listing style ────────────────────────────────────────────────
\lstdefinestyle{polya-python}{
  language=Python,
  basicstyle=\ttfamily\small,
  keywordstyle=\color{polyablue}\bfseries,
  commentstyle=\color{polyagray}\itshape,
  stringstyle=\color{polyagreen},
  numbers=left,
  numberstyle=\tiny\color{polyagray},
  numbersep=8pt,
  frame=single,
  framerule=0.4pt,
  rulecolor=\color{polyagray},
  breaklines=true,
  showstringspaces=false,
  tabsize=4,
}
\lstset{style=polya-python}
\title{Customer Segmentation via Clustering}
\author{uber-polya}
\date{2026-02-26}

\begin{document}

\maketitle
\tableofcontents
\newpage

% ── Phase A: Problem Formulation ─────────────────────────────────────
\section{Problem Formulation}

\begin{polyamodel}

\textbf{Problem Type}: Prove \hfill
\textbf{Domain}: Machine Learning


\end{polyamodel}

% ── Universe ──────────────────────────────────────────────────────────
\subsection{Universe}

\begin{itemize}
  \item $Elbow Analysis: \{k_range, inertias\}$
  \item $Kmeans: \{k, labels, centroids, silhouette, inertia\}$
  \item $Dbscan: \{eps, labels, n_clusters, n_noise, silhouette\}$
  \item $Gmm: \{n_components, labels, bic, aic, silhouette\}$
\end{itemize}

% ── Variables ─────────────────────────────────────────────────────────
\subsection{Variables}

\begin{longtable}{lll}
\toprule
\textbf{Symbol} & \textbf{Meaning} & \textbf{Type / Range} \\
\midrule
\endhead
$x$ & decision variable & $\mathbb{{R}}$ \\
\bottomrule
\end{longtable}

% ── Structure ─────────────────────────────────────────────────────────
\subsection{Structure}

- **Data**: \textasciitilde{}500 synthetic customers with 4 purchasing features
- **Features**: annual\_spending, purchase\_frequency, avg\_order\_value, days\_since\_last\_purchase
- **Natural segments**: 3 ground-truth clusters with overlapping boundaries
  - **High Value** (\textasciitilde{}150 customers): high spending, high frequency, low recency
  - **Moderate** (\textasciitilde{}200 customers): medium spending, medium frequency
  - **At Risk** (\textasciitilde{}150 customers): low spending, low frequency, high recency
- **Questions**: (1) How many segments exist? (2) Which method recovers them best? (3) What defines each segment?

% ── Mapping ───────────────────────────────────────────────────────────
\subsection{Real-World Mapping}

\begin{longtable}{ll}
\toprule
\textbf{Real-World Concept} & \textbf{Mathematical Object} \\
\midrule
\endhead
problem input & $instance parameters$ \\
\bottomrule
\end{longtable}

% ── Constraints ───────────────────────────────────────────────────────
\subsection{Constraints}

\begin{enumerate}
  \item $K-Means**: partition-based clustering minimizing within-cluster variance; assumes spherical, equal-size clusters$
  \hfill \textit{()}
  \item $DBSCAN**: density-based clustering that finds arbitrarily shaped clusters and identifies noise points$
  \hfill \textit{()}
  \item $GMM (Gaussian Mixture Model)**: probabilistic model assuming data is generated from a mixture of Gaussians$
  \hfill \textit{()}
  \item $Silhouette score**: measures how similar an object is to its own cluster vs. the nearest neighbor cluster (range -1 to +1)$
  \hfill \textit{()}
  \item $Elbow method**: heuristic for choosing k by finding the "knee" in the inertia curve$
  \hfill \textit{()}
  \item $Cluster profiles**: per-cluster summary statistics that translate numeric labels into business-meaningful segments$
  \hfill \textit{()}
\end{enumerate}

% ── Objective / Claim ─────────────────────────────────────────────────
\subsection{Claim}


% ── Approach ──────────────────────────────────────────────────────────
\subsection{Approach}

\begin{description}
  \item[Suggested approach:] K-Means + DBSCAN + GMM
  \item[Complexity class:] 
  \item[Available tools:] K-Means + DBSCAN + GMM
\end{description}
% ── Phase B: Solution ────────────────────────────────────────────────
\section{Solution}

\begin{polyaresult}

\textbf{Answer}: Solution computed


\textbf{Optimal}: No / Unknown \hfill
\textbf{Feasible}: No
\textbf{Algorithm}: K-Means + DBSCAN + GMM \quad $$

\textbf{Time}: 0.21494594000978395s


\end{polyaresult}

% ── Solution Details ──────────────────────────────────────────────────
\subsection{Solution Details}


Method: 1. StandardScaler: Zero-mean, unit-variance normalization of all features. Essential for distance-based methods (K-Means, DBSCAN) where feature scales differ by orders of magnitude.

2. Elbow method: Run K-Means for k=2..8, plot inertia (within-cluster sum of squares). The "elbow" where marginal improvement drops suggests the natural number of clusters.

3. K-Means (k=3): Lloyd's algorithm with k-means++ initialization. Fast, deterministic given seed, assumes spherical clusters. Reports centroids, labels, inertia, and silhouette score.

4. DBSCAN: Density-based clustering that discovers clusters of arbitrary shape. The eps parameter is set via k-distance heuristic (knee of the sorted 5-nearest-neighbor distance curve). Reports labels, number of clusters, number of noise points, and silhouette score.

5. GMM (3 components): Gaussian Mixture Model fit via Expectation-Maximization. Soft clustering that models each cluster as a multivariate Gaussian. Reports labels (hard assignment), BIC, AIC, and silhouette score.

6. Comparison: Methods ranked by silhouette score; best method selected.

7. Profiling: For the best method, compute mean feature values per cluster to build interpretable segment descriptions.

% ── Verification ──────────────────────────────────────────────────────
\subsection{Verification}

\begin{longtable}{lcl}
\toprule
\textbf{Check} & \textbf{Status} & \textbf{Detail} \\
\midrule
\endhead
Kmeans Silhouette Gt 04 & \passmark & Passed \\
Kmeans Found 3 Clusters & \passmark & Passed \\
Dbscan Found 2 To 5 Clusters & \passmark & Passed \\
Gmm Bic Finite & \passmark & Passed \\
Cluster Sizes Reasonable & \passmark & Passed \\
Profiles Distinguishable & \passmark & Passed \\
All Passed & \passmark & Passed \\
\bottomrule
\end{longtable}

% ── Phase C: Interpretation ──────────────────────────────────────────
\section{Interpretation}

% ── The Question & The Answer ─────────────────────────────────────────
\subsection{The Question}
- **Data**: \textasciitilde{}500 synthetic customers with 4 purchasing features
- **Features**: annual\_spending, purchase\_frequency, avg\_order\_value, days\_since\_last\_purchase
- **Natural segments**: 3 ground-truth clusters with overlapping boundaries
  - **High Value** (\textasciitilde{}150 customers): high spending, high frequency, low recency
  - **Moderate** (\textasciitilde{}200 customers): medium spending, medium frequency
  - **At Risk** (\textasciitilde{}150 customers): low spending, low frequency, high recency
- **Questions**: (1) How many segments exist? (2) Which method recovers them best? (3) What defines each segment?.

\subsection{The Answer}

\begin{polyainsight}[title={Bottom Line}]
A feasible solution was found.
\end{polyainsight}

% ── What This Means ───────────────────────────────────────────────────
\subsection{What This Means}

- Elbow analysis: inertia curve for k=2..8 shows clear elbow at k=3
- K-Means (k=3): silhouette score > 0.4, recovers 3 well-separated clusters
- DBSCAN: finds 2--5 clusters depending on eps (determined via k-distance heuristic)
- GMM (3 components): finite BIC/AIC, comparable silhouette to K-Means
- Best method: selected by highest silhouette score
- Cluster profiles: distinct mean feature values per cluster (distinguishable in standardized space)
- All 7 verification checks pass

1. StandardScaler: Zero-mean, unit-variance normalization of all features. Essential for distance-based methods (K-Means, DBSCAN) where feature scales differ by orders of magnitude.

2. Elbow method: Run K-Means for k=2..8, plot inertia (within-cluster sum of squares). The "elbow" where marginal improvement drops suggests the natural number of clusters.

3. K-Means (k=3): Lloyd's algorithm with k-means++ initialization. Fast, deterministic given seed, assumes spherical clusters. Reports centroids, labels, inertia, and silhouette score.

4. DBSCAN: Density-based clustering that discovers clusters of arbitrary shape. The eps parameter is set via k-distance heuristic (knee of the sorted 5-nearest-neighbor distance curve). Reports labels, number of clusters, number of noise points, and silhouette score.

5. GMM (3 components): Gaussian Mixture Model fit via Expectation-Maximization. Soft clustering that models each cluster as a multivariate Gaussian. Reports labels (hard assignment), BIC, AIC, and silhouette score.

6. Comparison: Methods ranked by silhouette score; best method selected.

7. Profiling: For the best method, compute mean feature values per cluster to build interpretable segment descriptions.

% ── Sensitivity Analysis ──────────────────────────────────────────────

% ── Figures ───────────────────────────────────────────────────────────

% ── Recommendations ───────────────────────────────────────────────────
\subsection{Recommendations}

\begin{enumerate}
  \item Review the model assumptions to ensure real-world fidelity.
\end{enumerate}

% ── Limitations ───────────────────────────────────────────────────────
\subsection{Limitations}

\begin{polyawarnbox}
\begin{itemize}
  \item Model assumes deterministic parameters; real-world variability not captured.
  \item Solution quality depends on the accuracy of input data.
\end{itemize}
\end{polyawarnbox}

% ── Appendix ─────────────────────────────────────────────────────────

\end{document}